% Complete proofs for independent families, product laws and mean-square bounds.
\section{Independence: proofs}
\label{sec:annex-p15-independent-random-variables}

Use the definitions of Chapter~\ref{sec:p15-independent-random-variables}.
The source independence development is Durrett's fixed author version
\cite[pp.~43--49]{DurrettPTEAuthor2019}. The proofs use the already
constructed probability measures and the full integration, product and
pi-lambda results; finite factorization is never treated as an unproved
existence theorem for a general infinite product.

\begin{annexitem}{annex:p15.independence-basic}{thm:p15-independence-basic}{Events, indicators, information and measurable null replacements}
\begin{proof}
Every event in \(\sigma(X_i)\) is an inverse image
\(X_i^{-1}(B_i)\) for a measurable value set \(B_i\), by the
definition and the proved inverse-image sigma-algebra theorem. Thus
the factorization required for a finite selection of events from
\(\sigma(X_i)\) is exactly the factorization in the definition
of independent random elements. This proves the equivalence for each
finite subfamily and hence for a countable family. If \(X_i\) is
\(\mathcal F_i\)-measurable, then \(\sigma(X_i)\subseteq\mathcal F_i\),
so independence of the larger sigma-algebras supplies that of the
smaller ones. Restricting indices or restricting each collection of
permitted events removes required tests and therefore preserves the
factorization property.

For complements, first work with a finite independent family of events
\(A_1,\ldots,A_n\) and complement just \(A_1\). A selected
intersection not using index one is unchanged. For a selected index
set \(J\) containing one, put
\(D=\bigcap_{j\in J\setminus\{1\}}A_j\). The empty case is
\(D=\Omega\). Finite measure subtraction and independence give
\[
 \begin{aligned}
 \mathbb P(A_1^c\cap D)
   &=\mathbb P(D)-\mathbb P(A_1\cap D)\\
   &=\prod_{j\in J\setminus\{1\}}\mathbb P(A_j)
       -\mathbb P(A_1)\prod_{j\in J\setminus\{1\}}\mathbb P(A_j)\\
   &=\mathbb P(A_1^c)
                      \prod_{j\in J\setminus\{1\}}\mathbb P(A_j).
 \end{aligned}
\]
This proves independence after one complementation. Relabeling the
indices and applying this argument successively handles any finite
selection of complements. In a countable family, each finite subfamily
involves only finitely many of the complemented coordinates, so the
same argument proves independence of the whole modified family.

For a Borel set \(B\), the inverse image
\(\{\mathbf1_{A_i}\in B\}\) is one of
\(\varnothing,A_i,A_i^c,\Omega\). In a finite intersection, an
empty factor makes the probability and the proposed product both zero.
Whole-space factors can be deleted. The other factors form an
intersection of events from a complemented independent subfamily, so
the preceding argument proves the required factorization. Thus the
indicators are independent. Conversely, choose the value set \(\{1\}\)
at every index of any finite subfamily of independent indicators to
recover the event factorization. Their generated sigma-algebras are
exactly \(\{\varnothing,A_i,A_i^c,\Omega\}=\sigma(\{A_i\})\),
with repetitions harmless when \(A_i\) is empty or full. The first
assertion now gives the sigma-algebra equivalence. This argument covers
the source's four complementary cases and every other indicator
inverse-image case.

An event independent of itself satisfies
\(\mathbb P(A)=\mathbb P(A)^2\). Its probability lies in
\([0,1]\), so it is zero or one. If \(\mathbb P(A)=0\),
monotonicity gives \(\mathbb P(A\cap B)=0\) for every \(B\).
If \(\mathbb P(A)=1\), its complement is null, and
\(\mathbb P(B)=\mathbb P(A\cap B)+\mathbb P(A^c\cap B)
=\mathbb P(A\cap B)\). These are precisely the independence
equalities in the two cases. The equivalence with conditional
probability for \(\mathbb P(B)>0\) follows by division by that
positive number, and its converse by multiplication. No conditional
value on a null event is assigned.

Finally suppose \(X_i=Y_i\) outside a measurable null set \(N_i\).
For any Borel \(B_i\), the inverse-image events differ only on
\(N_i\), so they have equal probabilities. For a finite index
set \(J\), the two corresponding intersection events differ only
on the measurable null union \(\bigcup_{i\in J}N_i\), and also
have equal probabilities. Factorization for the \(X_i\) therefore
gives factorization for the \(Y_i\). The one-coordinate assertion
gives equality of marginal laws. This proves the countable-family
conclusion through its finite subfamilies. The event definitions
are identical to the earlier finite ones; equivalence of the finite
random-variable point-mass criterion with all value-set tests was
already proved in Theorem~\ref{thm:p15-finite-independence-laws}.
\end{proof}
\end{annexitem}

\begin{annexitem}{annex:p15.independence-generators}{thm:p15-independence-generators}{Adjoining the whole space and extending one coordinate at a time}
\begin{proof}
Put \(\overline{\mathcal C}_i=\mathcal C_i\cup\{\Omega\}\).
For a finite selection \(A_i\in\overline{\mathcal C}_i\), retain
the index set \(K=\{i:A_i\ne\Omega\}\). Deleting the other
factors changes neither the intersection nor the probability product,
since \(\mathbb P(\Omega)=1\). Independence of the original
collections gives factorization over \(K\), also when \(K\)
is empty. Thus adjoining \(\Omega\) preserves independence.
Conversely restricting back to the original collections preserves it.
When \(\Omega\) is already present in each collection of a
finite family, fill every unused coordinate with \(\Omega\).
The complete-family equality then gives every subfamily equality,
establishing the first assertions.

Adjoining \(\Omega\) to a pi-system preserves its intersection
closure and does not change its generated sigma-algebra. Thus for
the finite theorem, the empty-family case is the identity one equals
one. For \(n\geq1\), suppose \(\mathcal C_1,\ldots,\mathcal C_n\)
all contain \(\Omega\). Fix \(A_i\in\mathcal C_i\) for
\(2\leq i\leq n\) and put \(D=\bigcap_{i=2}^n A_i\),
with \(D=\Omega\) when \(n=1\). Independence gives
\(\mathbb P(D)=\prod_{i=2}^n\mathbb P(A_i)\).
Define a class of events in the ambient sigma-algebra by
\[
 \mathcal L_D=\{A\in\mathcal F:
          \mathbb P(A\cap D)=\mathbb P(A)\mathbb P(D)\}.
\]
It contains \(\Omega\). For \(A\subseteq B\) in this class,
finite probability subtraction gives
\[
 \begin{aligned}
 \mathbb P((B\setminus A)\cap D)
   &=\mathbb P(B\cap D)-\mathbb P(A\cap D)\\
   &=(\mathbb P(B)-\mathbb P(A))\mathbb P(D)
     =\mathbb P(B\setminus A)\mathbb P(D).
 \end{aligned}
\]
For \(B_r\uparrow B\) in the class, continuity from below gives
\[
 \mathbb P(B\cap D)=\lim_r\mathbb P(B_r\cap D)
        =\lim_r\mathbb P(B_r)\mathbb P(D)
        =\mathbb P(B)\mathbb P(D).
\]
Thus \(\mathcal L_D\) is a lambda-system. Independence of the
original collections puts \(\mathcal C_1\) in this class.
The proved pi-lambda theorem gives
\(\sigma(\mathcal C_1)\subseteq\mathcal L_D\).

For every \(A_1\in\sigma(\mathcal C_1)\) and the fixed remaining
events, it follows that
\[
 \mathbb P\left(\bigcap_{i=1}^n A_i\right)
       =\mathbb P(A_1)\mathbb P(D)
       =\prod_{i=1}^n\mathbb P(A_i).
\]
The remaining events were arbitrary, and all the mixed collections
contain \(\Omega\), so the first part of the proof makes
\(\sigma(\mathcal C_1),\mathcal C_2,\ldots,\mathcal C_n\)
independent. Repeat the argument for the second coordinate, then the
third, and so on. At each step the collection being extended is a
pi-system; previously extended collections are sigma-algebras and
continue to contain \(\Omega\). After finitely many steps all
the generated sigma-algebras are independent. For a countable index
set, apply this finite result to each finite subset of indices.
This is exactly the definition of independence of the resulting family.
\end{proof}
\end{annexitem}

\begin{annexitem}{annex:p15.independence-thresholds}{thm:p15-independence-thresholds}{Generating the value sigma-algebras with threshold events}
\begin{proof}
Necessity follows directly from independence: each finite real
threshold is a Borel value set, and the infinite threshold is the
whole value space. For sufficiency using rational thresholds, let
\[
 \mathcal C_i=\{\{X_i\leq q\}:q\in\mathbb Q\}\cup\{\Omega\}.
\]
It is a pi-system, because the intersection of two finite-threshold
events uses their smaller rational threshold, and intersecting with
\(\Omega\) leaves the event unchanged. The hypothesis supplies
the complete-family factorization for every choice from these
collections. Since they contain \(\Omega\), it supplies all
subfamily factorizations as well. The generator theorem gives
independence of \(\sigma(\mathcal C_i)\).

The rational closed rays generate the Borel sigma-algebra on
\(\mathbb R\). The inverse-image generator result in
Theorem~\ref{thm:p15-measurability-criterion} therefore identifies
\(\sigma(\mathcal C_i)=\sigma(X_i)\). Independence of the
variables follows from Theorem~\ref{thm:p15-independence-basic}.
Testing every real threshold and infinity includes these rational
tests, so also suffices and retains the source's exact criterion.
The finite result applied to every finite list of distinct indices
proves the stated countable-family criterion.
\end{proof}
\end{annexitem}

\begin{annexitem}{annex:p15.independent-blocks}{thm:p15-independent-blocks}{Independent sigma-algebras from disjoint index blocks}
\begin{proof}
For each group \(I_j\), let \(\mathcal C_j\) be the collection
of finite intersections
\[
 \bigcap_{i\in K}A_i,
       \qquad K\subseteq I_j\text{ finite},\quad A_i\in\mathcal F_i.
\]
Include the empty intersection \(\Omega\). This is a pi-system.
To intersect two such expressions, take the union of their finite
index sets; at a repeated index intersect the two events in the
sigma-algebra \(\mathcal F_i\). Every member of
\(\bigcup_{i\in I_j}\mathcal F_i\) is one of these intersections,
and all these intersections belong to \(\mathcal G_j\). Hence
\(\sigma(\mathcal C_j)=\mathcal G_j\). If \(I_j\) is empty,
\(\mathcal C_j=\{\Omega\}\), whose generated sigma-algebra is
\(\{\varnothing,\Omega\}\).

Select finitely many different groups and one event \(C_j\in\mathcal C_j\)
from each, represented by its finite \(K_j\subseteq I_j\).
The index sets \(K_j\) are disjoint, and their union is finite.
Independence of the original sigma-algebras therefore gives
\[
 \begin{aligned}
 \mathbb P\left(\bigcap_j C_j\right)
   &=\prod_j\prod_{i\in K_j}\mathbb P(A_i)\\
   &=\prod_j\mathbb P\left(\bigcap_{i\in K_j}A_i\right)
     =\prod_j\mathbb P(C_j).
 \end{aligned}
\]
Both equalities include empty groups, whose contributions are one.
Thus the collections \(\mathcal C_j\) are independent. Applying
the generator theorem proves independence of the \(\mathcal G_j\).
If \(Y_j\) is \(\mathcal G_j\)-measurable, its generated
sigma-algebra is contained there, so the basic equivalence proves
independence of the \(Y_j\).

For a single-coordinate measurable transformation, composition makes
\(f_i(X_i)\) measurable with respect to \(\sigma(X_i)\).
The same equivalence gives the assertion. Directly from the definition,
for independent \(X,Y\) and measurable \(f,g\), any measurable
target sets \(B,C\) give
\[
 \begin{aligned}
 \mathbb P(f(X)\in B,g(Y)\in C)
   &=\mathbb P(X\in f^{-1}(B),Y\in g^{-1}(C))\\
   &=\mathbb P(X\in f^{-1}(B))\mathbb P(Y\in g^{-1}(C))\\
   &=\mathbb P(f(X)\in B)\mathbb P(g(Y)\in C).
 \end{aligned}
\]
The inverse-image sets are measurable by the hypotheses on \(f,g\),
so this is also the direct proof requested in source Exercise 2.1.6.
In the finite
real-block case, all coordinates of each block are measurable for
its \(\mathcal G_i\). The measurable-vector theorem makes the
block map into \(\mathbb R^{m(i)}\) measurable, and Borel
composition gives measurability of its image under \(f_i\).
The blocks are disjoint, so the preceding conclusion applies.

In the displayed past/future application, each partial sum up to
\(m\), their finite maximum, and its threshold indicator are
measurable functions of \(X_1,\ldots,X_m\). The increment uses
only \(X_{m+1},\ldots,X_n\). Their disjointness supplies the
required independence. The proof for countable groups still uses
only finite intersections and generated sigma-algebras. It does
not assert that every unspecified function of a countable block
is measurable; the theorem retains that premise for its output.
\end{proof}
\end{annexitem}

\begin{annexitem}{annex:p15.independent-product-law}{thm:p15-independent-product-law}{Identifying and constructing finite joint laws}
\begin{proof}
The map \(\omega\mapsto(X_1(\omega),\ldots,X_n(\omega))\)
is measurable for the finite product sigma-algebra: the inverse
image of a rectangle \(\prod_i B_i\) is the finite intersection
\(\bigcap_i\{X_i\in B_i\}\), and rectangles generate the
target domain. Thus its joint law is a probability measure there.
All marginal laws \(\mu_i\) are probability measures, so the
finite product measure has been constructed in
Theorem~\ref{thm:p15-finite-product-measures}.
If the variables are independent, on every measurable rectangle,
\[
 \begin{aligned}
 \mathbb P((X_1,\ldots,X_n)\in\textstyle\prod_iB_i)
   &=\prod_i\mathbb P(X_i\in B_i)\\
   &=\prod_i\mu_i(B_i)
     =(\mu_1\times\cdots\times\mu_n)(\textstyle\prod_iB_i).
 \end{aligned}
\]
The two probability measures agree on a generating pi-system.
The uniqueness part of Theorem~\ref{thm:p15-measure-uniqueness}
therefore makes them equal on every product measurable set.
Conversely, equality with the product law gives this rectangle
factorization. To extract any subfamily, take \(B_i=S_i\)
at the unused indices; those factors have measure one. This proves
independence in its full definition.

For finite construction, take the product space supplied by the
product-measure theorem and its coordinate maps. The coordinates
are measurable. Giving every coordinate other than \(i\) its
whole value space shows that the \(i\)-th law is \(\mu_i\).
The rectangle formula gives the required joint factorization, hence
independence. The constructed Stieltjes probabilities for admissible
\(F_i\) are a special case. Their bounded half-open interval
values are \(F_i(b_i)-F_i(a_i)\), so the product formula gives
the displayed finite-rectangle probability. This proves existence
for every finite list of such distributions.

For the countable-valued criterion, the finite product
\(S=\prod_{i=1}^n S_i\) is countable. Its product sigma-algebra
contains every singleton rectangle and therefore every subset, since
each subset is a countable union of singletons. The joint law and
the product of the marginals are two probability measures on this
domain. The assumed point-mass equality says that their values agree
on every singleton. Countable additivity, summing over the singletons
of any subset, proves equality on that subset. They are the same
joint law, and the converse just proved gives independence. This
also proves the claimed agreement with the finite-valued criterion.
\end{proof}
\end{annexitem}

\begin{annexitem}{annex:p15.independent-expectations}{thm:p15-independent-expectations}{The product-law integral and all finite product cases}
\begin{proof}
The joint law of \((X,Y)\) is \(\mu\times\nu\) by the
preceding theorem. The law-integration identity in
Theorem~\ref{thm:p15-expectation-law-integral} therefore gives
\[
 \mathbb E[h(X,Y)]=\int h\,d(\mu\times\nu)
\]
in the nonnegative or integrable case. Apply Tonelli in the first
case and Fubini in the second. The factors are finite probability
measures, hence sigma-finite. These theorems give both iterated
orders, their measurability, and the exact a.e. section conventions
for the signed case.

Suppose \(f,g\) are nonnegative real Borel functions. Put
\(A=\int f\,d\mu\), \(B=\int g\,d\nu\), with values in
\([0,\infty]\). For every fixed \(y\), the finite constant
\(g(y)\) can be taken outside the nonnegative inner integral:
\[
 \int f(x)g(y)\,d\mu(x)=A g(y).
\]
If \(A<\infty\), nonnegative scaling in the outer integral gives
\(\int A g\,d\nu=AB\). If \(A=\infty\), the measurable
functions \(r g\), \(r=1,2,\ldots\), increase to the
specified \(A g\). Monotone convergence gives
\(\int A g\,d\nu=\lim_r rB\), which is zero for \(B=0\)
and infinity otherwise. It is exactly \(AB\) under the adopted
zero-times-infinity convention. Together with the law formula,
this proves \(\mathbb E[f(X)g(Y)]=\mathbb E[f(X)]\mathbb E[g(Y)]\)
in every nonnegative case.

For signed real \(f,g\) with the stated absolute moments, apply
this result to \(|f|,|g|\) to obtain
\[
 \mathbb E|f(X)g(Y)|=
             \mathbb E|f(X)|\,\mathbb E|g(Y)|<\infty.
\]
Thus the signed product is integrable before Fubini is used on it.
Its inner integral is \(g(y)\int f\,d\mu\), and its outer
integral is \((\int f\,d\mu)(\int g\,d\nu)\), by signed
linearity. The law formula for each factor gives the stated identity.
All these signed numbers are finite.

For finitely many independent nonnegative real variables, induct on
their number. The assertion for one variable is tautological. The
disjoint-block theorem makes \(X_1\) independent of
\(Y=X_2\cdots X_n\); this is a real measurable nonnegative
function of the remaining finite block. The two-factor result and
the induction hypothesis give
\[
 \mathbb E[X_1\cdots X_n]=\mathbb E[X_1]\mathbb E[Y]
                      =\prod_{i=1}^n\mathbb E[X_i].
\]
Finite nonnegative extended multiplication is associative with the
declared convention: any zero factor makes the result zero; in its
absence a positive infinite factor gives infinity; otherwise it is
ordinary positive multiplication. This includes all infinite-mean cases.

For independent individually integrable real variables, their absolute
values remain independent by measurable transformation. The result
just proved gives
\(\mathbb E\prod_i|X_i|=\prod_i\mathbb E|X_i|<\infty\),
and the same holds for each subproduct. Consequently
\(Y=X_2\cdots X_n\) is integrable. Use the signed two-factor
identity for \(X_1,Y\), then induction on the remaining signed
subproduct, obtaining the desired signed factorization.

For the converse, choose bounded Borel functions
\(f=\mathbf1_B\), \(g=\mathbf1_C\) for arbitrary Borel
\(B,C\). The assumed expectation identity becomes
\(\mathbb P(X\in B,Y\in C)=\mathbb P(X\in B)\mathbb P(Y\in C)\),
which is independence. For a countable family, apply the same
indicator choice at every index of any finite subfamily. Conversely,
independence gives independence of its bounded Borel transforms;
they are integrable on a probability space, so the finite product
formula proves every bounded-function test. This proves the last
characterization in both directions.
\end{proof}
\end{annexitem}

\begin{annexitem}{annex:p15.independent-moments}{thm:p15-independent-moments}{Finite second moments, covariance and the square expansion}
\begin{proof}
Finite second moments imply finite first absolute moments by
Cauchy--Schwarz with the constant one, as already established in
Theorem~\ref{thm:p15-expectation-bounds}. The same inequality for
\(X,Y\) gives \(\mathbb E|XY|<\infty\).
Put \(\mu_X=\mathbb E X\), \(\mu_Y=\mathbb E Y\).
The centered variables have finite second moments, equal to their
finite variances. Cauchy--Schwarz for these centered variables gives
\[
 \mathbb E|(X-\mu_X)(Y-\mu_Y)|
       \leq\sqrt{\operatorname{Var}(X)\operatorname{Var}(Y)}.
\]
Thus covariance is defined and finite. Taking the absolute value
of its integral gives the stated bound. Expanding its integrable
product and using signed linearity yields
\[
 \mathbb E[(X-\mu_X)(Y-\mu_Y)]
    =\mathbb E[XY]-\mu_X\mathbb E Y-\mu_Y\mathbb E X+\mu_X\mu_Y
    =\mathbb E[XY]-\mu_X\mu_Y.
\]
Division by the positive variance product square root proves the
correlation bound when it is defined. For independent \(X,Y\),
the preceding expectation theorem gives
\(\mathbb E[XY]=\mathbb E X\mathbb E Y\), so covariance is zero.

For the finite weighted sum, the elementary identity
\[
 n\sum_{i=1}^n v_i^2-\left(\sum_{i=1}^n v_i\right)^2
                         =\sum_{i<j}(v_i-v_j)^2\geq0
\]
shows, on putting \(v_i=a_iX_i\), that its square is bounded
by an integrable finite sum \(n\sum_i a_i^2X_i^2\).
Thus its second moment is finite. Its mean is \(\sum_i a_i\mu_i\),
where \(\mu_i=\mathbb E X_i\), by finite signed linearity.
Expanding its centered square gives
\[
 \left(\sum_i a_i(X_i-\mu_i)\right)^2
   =\sum_i a_i^2(X_i-\mu_i)^2
       +2\sum_{i<j}a_ia_j(X_i-\mu_i)(X_j-\mu_j).
\]
Each squared term is integrable, and each cross term is integrable
by the centered Cauchy--Schwarz argument above. Integrating this
finite identity therefore gives exactly the variance formula.
If every distinct pair is uncorrelated, each cross covariance is
zero. Pairwise independence supplies this condition by the two-variable
result and requires no stronger mutual-independence premise.
\end{proof}
\end{annexitem}

\begin{annexitem}{annex:p15.independence-counterexamples}{thm:p15-independence-counterexamples}{The complete finite countermodels}
\begin{proof}
The uniform measure on the eight binary triples is the finite
product of three fair binary probabilities. Its coordinate variables
are independent by the finite product construction. Each event
equating two coordinates has four outcomes, since their common bit
and the remaining bit can be chosen in two ways each. Its probability
is \(4/8=1/2\). Requiring any two of \(A_1,A_2,A_3\) forces
all three coordinates to agree, and conversely agreement satisfies
all three events. Exactly the outcomes \((0,0,0)\), \((1,1,1)\)
remain. Thus every pairwise intersection and the triple intersection
have probability \(2/8=1/4\). This equals the product of two
marginals but differs from the product of three. The indicator
equivalence transfers precisely those pairwise and mutual conclusions
to the three indicators.

For the four-outcome generator example, each of \(A,B,C\) has
probability one half, and \(A\cap C=B\cap C=\{1\}\) has
probability one quarter. Hence every choice of one event from
\(\{A,B\}\) and one from \(\{C\}\) factors correctly;
single-collection and empty-subfamily tests are automatic. These
two collections are independent as defined. But
\(A\cap B=\{1\}\) belongs to \(\sigma(\{A,B\})\), and
\(C\in\sigma(\{C\})\), while
\[
 \mathbb P(\{1\}\cap C)=1/4
       \ne(1/4)(1/2)=\mathbb P(\{1\})\mathbb P(C).
\]
Their generated sigma-algebras are not independent. The intersection
\(A\cap B\) is absent from \(\{A,B\}\), exhibiting exactly
the failed pi-system hypothesis. Adding \(\Omega\) to the two
collections does not change these facts.

For the third example, the five possibly nonzero table entries are
nonnegative and total \(2a+2b+c=1\), so they define a probability
measure on the nine displayed value pairs. Its coordinate variables
have finite moments because they are bounded by one in absolute value.
Every nonzero mass has either \(X=0\) or \(Y=0\), hence
\(XY=0\) with probability one. The marginal masses of \(X\)
at \(1,-1\) are both \(a\), and those of \(Y\) at
\(1,-1\) are both \(b\). Therefore
\[
 \mathbb E X=a-a=0,\quad \mathbb E Y=b-b=0,\quad
 \mathbb E X^2=2a,\quad \mathbb E Y^2=2b.
\]
The asserted variances, covariance and correlation follow from the
moment theorem; both variances are positive by \(a,b>0\).
Yet the joint point mass at \((1,1)\) is zero, whereas the
product of its marginal masses is \(ab>0\). This violates the
necessary independence factorization. Zero covariance, even with
positive finite variances and a well-defined zero correlation, does
not imply independence.

Finally, a fair Bernoulli \(Z\) exists on a two-point uniform
space, or as any coordinate in the first example. Its event
\(A=\{Z=1\}\) has probability one half. The equality required
for independence of \(Z\) with itself would give
\(\mathbb P(A)=\mathbb P(A)^2\), but \(1/2\ne1/4\).
Thus identical outputs from overlapping source groups need not be
independent. This is a failure of the required disjointness, not of
measurability of either output.
\end{proof}
\end{annexitem}

\begin{annexitem}{annex:p15.dyadic-independent-coins}{thm:p15-dyadic-independent-coins}{Dyadic cells specify all finite bit patterns}
\begin{proof}
Put \(q_n(u)=\lfloor2^nu\rfloor\) for \(n\geq0\), so
\(q_0(u)=0\) on \((0,1)\). Each \(q_n\) is measurable:
the inverse image of any one of its finitely many values is a
half-open dyadic interval intersected with \((0,1)\).
Writing \(2^{n-1}u=q_{n-1}(u)+r\), with \(0\leq r<1\),
gives
\[
 q_n(u)=2q_{n-1}(u)+\lfloor2r\rfloor.
\]
The last floor is either zero or one. Thus
\(Y_n=q_n-2q_{n-1}\) is a measurable binary variable and has
the same parity as \(q_n\), proving agreement with the source
definition at every point.

Fix a positive integer \(m\). Iterating this recurrence, starting
from \(q_0=0\), yields
\[
 q_m(u)=\sum_{j=1}^m2^{m-j}Y_j(u).
\]
The map from \(m\) binary digits to this integer is a bijection
onto \(\{0,\ldots,2^m-1\}\). One proof successively takes
the integer's remainder upon division by two to recover its last
digit, and its quotient to recover the preceding digits; induction
gives both existence and uniqueness. Consequently the prefix event
with specified digits \(\varepsilon_1,\ldots,\varepsilon_m\)
is exactly
\[
 C_k=\left[\frac{k}{2^m},\frac{k+1}{2^m}\right)\cap(0,1),
 \qquad k=\sum_{j=1}^m2^{m-j}\varepsilon_j.
\]
Indeed, membership in this cell is \(q_m=k\); the displayed
recurrence and uniqueness of the digit representation are equivalent
to all the prescribed prefix digits. The cells are disjoint and
partition \((0,1)\). Their probabilities are all \(2^{-m}\)
by the constructed interval length measure. Removing zero or one
at the outer boundaries does not change their lengths.

For nonempty finite \(J\subseteq\mathbb N_+\), choose
\(m=\max J\). Prescribing its \(\#J\) bits leaves
\(m-\#J\) prefix digits free. Exactly \(2^{m-\#J}\)
of the disjoint prefix cells therefore realize the prescribed
pattern. Finite additivity gives its probability
\(2^{m-\#J}2^{-m}=2^{-\#J}\). For an empty \(J\),
the event is the whole space and the same formula is one.
In particular every individual bit has both values with probability
one half.

To obtain independence of random variables, consider arbitrary
Borel value sets at finitely many indices. For each binary variable
the corresponding event is empty, the whole space, or one specified
bit value. An empty event gives zero on both sides of the
factorization. Whole-space events can be deleted. The remaining
specified-bit event has exactly the product probability just proved.
This verifies all finite subfamily tests, hence mutual independence
of the countable family. Every variable was defined on the existing
unit interval; no infinite-product extension theorem has been assumed.
The half-open cells and the floor formula consistently assign the
terminating binary choice at dyadic points, so no pointwise definition
is missing there.
\end{proof}
\end{annexitem}

\begin{annexitem}{annex:p15.independent-l2-weak-law}{thm:p15-independent-l2-weak-law}{A moment bound and the variance calculation for averages}
\begin{proof}
For real \(Z_n\) and \(p>0\), the function \(|Z_n|^p\)
is nonnegative and measurable. Markov's inequality at the positive
threshold \(\varepsilon^p\) gives
\[
 \mathbb P(|Z_n|\geq\varepsilon)
             \leq\varepsilon^{-p}\mathbb E|Z_n|^p.
\]
Under the stated convergence of moments, the right side tends to
zero for every \(\varepsilon>0\). This implies the required
convergence in probability, including its strict-threshold definition.
This is the source's Lemma 2.2.2 with the power applied to the
nonnegative absolute value.

For the sequence in the second assertion, finite linearity gives
\(\mathbb E S_n=n\mu\). The covariance formula for a finite
sum and pairwise zero covariances give
\(\operatorname{Var}(S_n)=\sum_{i=1}^n\operatorname{Var}(X_i)\leq nC\).
Scaling the variance and centering at the actual mean yield
\[
 \mathbb E[(S_n/n-\mu)^2]
       =\operatorname{Var}(S_n/n)
       =n^{-2}\operatorname{Var}(S_n)\leq C/n.
\]
This proves mean-square convergence. Applying the first Markov
calculation with \(p=2\) gives the displayed probability bound
and convergence in probability. Only the specified pairwise
uncorrelated second moments were used; no stronger independence
assumption enters the proof.

For the constructed bits, \(Y_i^2=Y_i\), so their mean is
one half and their variance is \(1/2-(1/2)^2=1/4\).
Their independence and finite second moments make them pairwise
uncorrelated by the preceding moment theorem. Substitution of
\(\mu=1/2\), \(C=1/4\) gives the final empirical-proportion
bound exactly.
\end{proof}
\end{annexitem}
